%!TEX root = ../LLMSeminar.tex
% ──────────────────────────────────────────────────────────────
%  Section 6 — Part II opener: Recap and the Question of Agency
% ──────────────────────────────────────────────────────────────

\bigskip
\begin{center}
  {\Large\sffamily\bfseries\color{cgblue} Part II \;---\; Anatomy of an Agent}\\[0.4em]
  \textit{(Second seminar in the series)}
\end{center}
\bigskip

\section{Recap and the question of agency}\label{sec:recap-agency}

\subsection{Where we left off}

Before we proceed it is worth stating, once more, the picture we
arrived at in the first seminar (\S\ref{sec:what-is-llm}--%
\S\ref{sec:chat-illusion}).

\begin{greybox}[Recap of Seminar~I]
  An LLM is a stateless, nondeterministic function
  $f : \Str \to \Str$.  It has no memory of any prior call and
  responds by sampling from a learned distribution over the next token.
  The phenomenon of ``chat'' is an \emph{illusion} maintained on the
  client side: every turn, the client re-sends the entire conversation
  history---the \emph{context window}---back to the function.  The
  function never remembers; the client remembers \emph{for} it.
\end{greybox}

The mental model worth holding onto is---\emph{highly motivated but
not trusted}.  The system you are talking to has read the entire
internet and will, with great confidence and articulacy, attempt
anything you ask of it; whether the answer is correct is a separate
question.

\subsection{From oracle to actor}

So far the LLM is an oracle: you supply context, it emits text.  The
\emph{only} thing the function can do, in a literal sense, is put
words on your screen.  This is the agency it has when used as a
plain chatbot:

\begin{center}
\begin{tikzpicture}[node distance=1.4cm]
  \node[sysbox fill] (chat) {Chat UI};
  \node[sysbox fill, right=4cm of chat] (llm) {LLM \;{\scriptsize $f(s)\to s$}};
  \draw[sysarrow] (chat.east) -- (llm.west)
    node[midway, above, font=\scriptsize] {prompt};
  \draw[sysarrow] (llm.west) ++(0,-0.55) -- ($(chat.east)+(0,-0.55)$)
    node[midway, below, font=\scriptsize] {tokens};
\end{tikzpicture}
\end{center}

We now want to build a programme that turns the function \emph{into
something that does things}---reads files, writes files, runs code,
queries the network.  The great surprise of 2024--2025 is that this
extension is essentially trivial in concept and that, once one
crosses this conceptual line, an enormous range of behaviour
becomes accessible.

\begin{definition}[Agency, in the present sense]\label{def:agency}
  We grant an LLM \emph{agency} when we wrap its single-call
  interface in a programme that interprets some of the function's
  output as \emph{requests for action}, performs those actions on
  the LLM's behalf, and feeds the results back into the next call.
  This wrapping programme is variously called the \emph{harness},
  the \emph{client}, or the \emph{scaffolding}.  Every coding agent
  in production today is exactly this: a harness around the function~$f$.
\end{definition}

\begin{keyresult}[Two facts about agency]
  \begin{enumerate}[(i)]
    \item The LLM never \emph{itself} executes anything.  It only
      emits tokens.  All effects on the outside world are produced
      by the harness, by virtue of the harness choosing to
      interpret some of the LLM's tokens as a request to act.
    \item Whether a particular LLM is \emph{good} at this depends
      on whether it has been \emph{trained} to.  Older models---
      \texttt{GPT-4o} of mid-2024, for example---understand
      English perfectly well but have not been taught a stable
      vocabulary for tool requests; modern models have, and the
      vendors have largely converged on a shared schema
      (cf.~\S\ref{sec:agent-loop}).
  \end{enumerate}
\end{keyresult}

\begin{intuition}[The harness is the hands]
  Think of the LLM as a brilliant but disembodied advisor on the
  far end of a phone line.  It can talk you through what to do
  but cannot, by itself, lift a finger.  The harness is the body:
  it has hands, eyes, a network connection, a file system.  When
  the advisor says ``read this file,'' the body reads it; when it
  says ``run this code,'' the body runs it; and the result is
  reported back over the phone line.
\end{intuition}

\subsection{Today's goal}

By the end of this lecture the following should be true.

\begin{enumerate}[(1)]
  \item You can describe the architecture of a coding agent
    in pseudocode---all of it---without omitting anything
    essential.
  \item You can write your own, in any language, in roughly
    a hundred lines.
  \item You have begun the deeper work of \emph{calibrating
    your expectations}: developing personal benchmarks that
    track what these systems can and cannot do, so that you
    are neither caught off guard nor swept along by hype.
  \item You begin to think about the science you might
    attempt with such tools, particularly at scales of
    ambition that were impossible for small groups before.
\end{enumerate}

\noindent The third and fourth points are, in my view, where
the real work of the seminar lies.  The first two are
prerequisites.
